% VSIM_EXAMPLES.tex
% VSEPR-SIM v5.14.1  |  VSIM Scripting Language — Examples
% Companion to VSIM_EXAMPLES.md
%
% Build:  pdflatex VSIM_EXAMPLES.tex  (twice for TOC/hyperlinks)
% Deps:   tcolorbox, listings, xcolor, geometry, hyperref, booktabs,
%         fontenc, inputenc, microtype, parskip, enumitem, mdframed
%
\documentclass[11pt,a4paper]{article}

%──────────────────────────────────────────────────────────────────────────────
% Packages
%──────────────────────────────────────────────────────────────────────────────
\usepackage[T1]{fontenc}
\usepackage[utf8]{inputenc}
\usepackage[margin=2.2cm]{geometry}
\usepackage{microtype}
\usepackage{parskip}
\usepackage{xcolor}
\usepackage{booktabs}
\usepackage{enumitem}
\usepackage{mdframed}
\usepackage{listings}
\usepackage[most]{tcolorbox}
\usepackage[colorlinks=true,linkcolor=vsimblue,urlcolor=vsimblue,
			citecolor=vsimblue]{hyperref}

%──────────────────────────────────────────────────────────────────────────────
% Colour palette
%──────────────────────────────────────────────────────────────────────────────
\definecolor{vsimblue}   {HTML}{1A73E8}
\definecolor{vsimgreen}  {HTML}{34A853}
\definecolor{vsimorange} {HTML}{FA7B17}
\definecolor{vsimgrey}   {HTML}{5F6368}
\definecolor{vsimltgrey} {HTML}{F1F3F4}
\definecolor{vsimcode}   {HTML}{202124}
\definecolor{vsimcodebg} {HTML}{F8F9FA}
\definecolor{vsimrule}   {HTML}{DADCE0}
\definecolor{vsimkw}     {HTML}{1967D2}
\definecolor{vsimstr}    {HTML}{188038}
\definecolor{vsimcmt}    {HTML}{80868B}
\definecolor{vsimnum}    {HTML}{D93025}
\definecolor{vsimwarn}   {HTML}{F29900}
\definecolor{vsimdoctrinebg}{HTML}{FFF8E1}
\definecolor{vsimdoctrinetop}{HTML}{F29900}

%──────────────────────────────────────────────────────────────────────────────
% VSIM listings style
%──────────────────────────────────────────────────────────────────────────────
\lstdefinelanguage{vsim}{
  sensitive    = true,
  morecomment  = [l]{\#},
  morestring   = [b]",
  morekeywords = {
	project, material, run, simulation, molecule,
	cell, boundary, pbc, environment, excite, chemistry, system,
	observe, variance, while, batch, job,
	export, visual, report, kernel, trace,
	analysis, structure, sampling, scale_sampling, inference,
	ikk_end_tag, ivec,
	object, macro, chemical, fundamental,
	carrier, action, information,
	override, particle, raw,
	room, N_evolution, dissolution
  },
}

\lstdefinestyle{vsimlst}{
  language         = vsim,
  basicstyle       = \ttfamily\small\color{vsimcode},
  backgroundcolor  = \color{vsimcodebg},
  commentstyle     = \color{vsimcmt}\itshape,
  keywordstyle     = \color{vsimkw}\bfseries,
  stringstyle      = \color{vsimstr},
  numberstyle      = \color{vsimnum},
  showstringspaces = false,
  breaklines       = true,
  breakatwhitespace= true,
  tabsize          = 4,
  xleftmargin      = 4pt,
  xrightmargin     = 4pt,
  frame            = none,
  aboveskip        = 0pt,
  belowskip        = 0pt,
}

\lstdefinestyle{terminallst}{
  basicstyle       = \ttfamily\footnotesize\color{white},
  backgroundcolor  = \color{vsimcode},
  showstringspaces = false,
  breaklines       = true,
  frame            = none,
  aboveskip        = 0pt,
  belowskip        = 0pt,
  xleftmargin      = 6pt,
  xrightmargin     = 6pt,
}

\lstdefinestyle{jsonlst}{
  basicstyle       = \ttfamily\footnotesize\color{vsimcode},
  backgroundcolor  = \color{vsimcodebg},
  showstringspaces = false,
  breaklines       = true,
  frame            = none,
  aboveskip        = 0pt,
  belowskip        = 0pt,
  xleftmargin      = 4pt,
  xrightmargin     = 4pt,
  morestring       = [b]",
}

%──────────────────────────────────────────────────────────────────────────────
% tcolorbox styles
%──────────────────────────────────────────────────────────────────────────────

% Main script box
\tcbset{
  scriptbox/.style={
	enhanced, breakable,
	colback   = vsimcodebg,
	colframe  = vsimrule,
	arc       = 3pt,
	boxrule   = 0.6pt,
	left      = 4pt, right = 4pt, top = 4pt, bottom = 4pt,
	before upper = {\vspace{-2pt}},
  },
  % Header strip for the script box
  scriptheader/.style={
	attach boxed title to top left={yshift=-2pt, xshift=6pt},
	boxed title style={
	  colback=vsimblue, colframe=vsimblue,
	  arc=2pt, boxrule=0pt,
	  left=4pt, right=4pt, top=2pt, bottom=2pt,
	},
	fonttitle=\ttfamily\bfseries\small\color{white},
  },
}

% Terminal output box
\tcbset{
  termbox/.style={
	enhanced, breakable,
	colback   = vsimcode,
	colframe  = vsimcode,
	arc       = 3pt,
	boxrule   = 0pt,
	left      = 6pt, right = 6pt, top = 4pt, bottom = 4pt,
  },
}

% Doctrine / callout box
\tcbset{
  doctrinebox/.style={
	enhanced, breakable,
	colback   = vsimdoctrinebg,
	colframe  = vsimdoctrinetop,
	arc       = 3pt,
	boxrule   = 1pt,
	left      = 8pt, right = 8pt, top = 4pt, bottom = 4pt,
	fontupper = \small,
  },
}

% Key-fields table box
\tcbset{
  tablebox/.style={
	enhanced, breakable,
	colback   = vsimltgrey,
	colframe  = vsimrule,
	arc       = 3pt,
	boxrule   = 0.4pt,
	left      = 6pt, right = 6pt, top = 4pt, bottom = 4pt,
  },
}

% Example header banner
\tcbset{
  exampleheader/.style={
	enhanced,
	colback   = vsimblue,
	colframe  = vsimblue,
	arc       = 4pt,
	boxrule   = 0pt,
	left      = 10pt, right = 10pt, top = 6pt, bottom = 6pt,
	fontupper = \bfseries\large\color{white},
  },
}

% Variation / tip box
\tcbset{
  tipbox/.style={
	enhanced,
	colback   = {green!7!white},
	colframe  = vsimgreen,
	arc       = 3pt,
	boxrule   = 0.8pt,
	left      = 8pt, right = 8pt, top = 3pt, bottom = 3pt,
	fontupper = \small,
  },
}

%──────────────────────────────────────────────────────────────────────────────
% Helpers
%──────────────────────────────────────────────────────────────────────────────
\newcommand{\exsection}[3]{%
  % #1 = label  #2 = number  #3 = title
  \newpage
  \phantomsection
  \addcontentsline{toc}{subsection}{#2 — #3}%
  \begin{tcolorbox}[exampleheader]
	#2 \quad #3
  \end{tcolorbox}%
}

\newcommand{\vsimscript}[2]{%
  % #1 = filename  #2 = caption
  \begin{tcolorbox}[scriptbox, scriptheader, title={#1}]
	\lstinputlisting[style=vsimlst]{#2}
  \end{tcolorbox}%
}

% Inline script using lstlisting directly (for embedded code)
\newenvironment{vsimblock}[1][script.vsim]{%
  \begin{tcolorbox}[scriptbox, scriptheader, title={#1}]%
  \lstset{style=vsimlst}%
  \lstset{aboveskip=0pt,belowskip=0pt}%
}{%
  \end{tcolorbox}%
}

\newenvironment{termblock}{%
  \begin{tcolorbox}[termbox]%
  \lstset{style=terminallst}%
}{%
  \end{tcolorbox}%
}

\newenvironment{doctrine}{%
  \begin{tcolorbox}[doctrinebox]%
  \textbf{\color{vsimdoctrinetop}Doctrine\quad}%
}{%
  \end{tcolorbox}%
}

\newenvironment{tip}{%
  \begin{tcolorbox}[tipbox]%
  \textbf{\color{vsimgreen}Variation / Tip\quad}%
}{%
  \end{tcolorbox}%
}

% Key-fields table (tabular wrapper)
\newenvironment{keyfields}{%
  \begin{tcolorbox}[tablebox]%
  \footnotesize%
  \begin{tabular}{@{}lp{4cm}p{7.5cm}@{}}%
  \toprule%
  \textbf{Field} & \textbf{Value used} & \textbf{Effect} \\%
  \midrule%
}{%
  \bottomrule%
  \end{tabular}%
  \end{tcolorbox}%
}

\newenvironment{keyfieldsnv}{%
  \begin{tcolorbox}[tablebox]%
  \footnotesize%
  \begin{tabular}{@{}lp{9cm}@{}}%
  \toprule%
  \textbf{Field} & \textbf{Effect} \\%
  \midrule%
}{%
  \bottomrule%
  \end{tabular}%
  \end{tcolorbox}%
}

%──────────────────────────────────────────────────────────────────────────────
% Title & document
%──────────────────────────────────────────────────────────────────────────────
\title{%
  {\Huge\bfseries\color{vsimblue} VSIM Scripting Language}\\[6pt]
  {\large Examples — Purpose-Driven Reference}\\[3pt]
	{\normalsize VSEPR-SIM v5.14.1 \quad|\quad companion: \texttt{VSIM\_EXAMPLES.md}}%
}
\author{VSEPR-SIM Development}
\date{\today}

\begin{document}
\maketitle
\thispagestyle{empty}

\begin{mdframed}[linecolor=vsimrule, backgroundcolor=vsimltgrey,
				 skipabove=4pt, skipbelow=4pt]
\textbf{How to read this document.}
Each example is self-contained and can be saved to a \texttt{.vsim} file and
run with \texttt{vsper run <script>} or validated with
\texttt{vsper validate <script>}.
Jump directly to the section you need — no serial reading required.
\end{mdframed}

\tableofcontents
\clearpage

%══════════════════════════════════════════════════════════════════════════════
\section*{Contents overview}
\addcontentsline{toc}{section}{Contents overview}
%══════════════════════════════════════════════════════════════════════════════

\begin{center}
\footnotesize
\begin{tabular}{@{}clp{8cm}@{}}
\toprule
\textbf{\#} & \textbf{Title} & \textbf{Concepts covered} \\
\midrule
E01 & Hello Atom              & Absolute minimum · static GL spin viewer \\
E02 & H\textsubscript{2}O Molecule Relax
							  & FIRE solver · terminal chart · per-atom override \\
E03 & NaCl Crystal --- Level-0 Intent
							  & Registry resolver · PBC · \texttt{research\_report} \\
E04 & Silicon Diamond with RDF
							  & Covalent prototype · Tersoff · \texttt{[observe]} \\
E05 & Argon Gas NVT           & Thermostat · MSD · \texttt{[[simulation.molecule]]} \\
E06 & TiO\textsubscript{2} Oxidation + While Loop
							  & \texttt{[chemistry]} · \texttt{[variance]} · \texttt{[while]} \\
E07 & Batch Temperature Sweep & \texttt{[batch]} · sweep axis · aggregate report \\
E08 & GL Visualizer Controls  & \texttt{gl\_interactive} · spin · orbit · overlay \\
E09 & IKK Identity Vector     & \texttt{[analysis.ivec]} · \texttt{.identity.json} · WO-75B \\
E10 & MCF-CAI Object State    & \texttt{[object.<layer>.<basis>]} · sidecar doctrine · WO-76 \\
E11 & Full Research Pipeline  & All major sections combined end-to-end \\
QR  & Quick-Reference Card    & One-page field cheatsheet \\
\bottomrule
\end{tabular}
\end{center}

%══════════════════════════════════════════════════════════════════════════════
\exsection{e01}{E01}{Hello Atom}
%══════════════════════════════════════════════════════════════════════════════

\textbf{Goal:} the absolute minimum script that opens a GL window with a
spinning hydrogen atom.

\textit{Key concepts:} \texttt{[project]} · \texttt{[material]} · \texttt{[run]}
· \texttt{[visual]} · \texttt{gl\_spin}

\begin{tcolorbox}[scriptbox, scriptheader, title={e01\_hello\_atom.vsim}]
\begin{lstlisting}[style=vsimlst]
# e01_hello_atom.vsim — minimum viable VSIM script
[project]
name    = "e01_hello_atom"
version = "v5.14.1"

[material]
formula   = "H"
prototype = "noble_gas"
phase     = "gas"

[run]
mode      = "md"
max_steps = 1
dt_fs     = 1.0
converge  = false    # no convergence needed — display run only

[[simulation.molecule]]
formula     = "H"
count       = 1
temperature = 0.0
lattice     = "none"

[export]
write_xyz  = true
output_dir = "out/e01"

[visual]
output_type       = "gl_live_60fps"
gl_spin           = true
gl_spin_axis      = "y"
gl_spin_deg_per_s = 45.0
gl_show_axes      = true
gl_window_width   = 960
gl_window_height  = 720
\end{lstlisting}
\end{tcolorbox}

\begin{keyfields}
\texttt{mode} + \texttt{max\_steps=1} & & One step; no integration \\
\texttt{converge = false}             & & Never exits on convergence \\
\texttt{output\_type = "gl\_live\_60fps"} & & Smooth 60\,fps OpenGL window \\
\texttt{gl\_spin = true}              & & Continuous viewer rotation \\
\texttt{gl\_spin\_deg\_per\_s}        & \texttt{45.0} & Try \texttt{180.0} for fast spin \\
\end{keyfields}

\begin{tip}
Change \texttt{gl\_spin\_axis} to \texttt{"x"} or \texttt{"z"}, or set
\texttt{gl\_spin\_deg\_per\_s = -30.0} for reverse spin.
\end{tip}

%══════════════════════════════════════════════════════════════════════════════
\exsection{e02}{E02}{H\textsubscript{2}O Molecule Relax}
%══════════════════════════════════════════════════════════════════════════════

\textbf{Goal:} relax a water molecule to its energy minimum and display live
convergence in the terminal.

\textit{Key concepts:} FIRE solver · \texttt{terminal\_chart}
· \texttt{[[override.particle]]} · export flags

\begin{tcolorbox}[scriptbox, scriptheader, title={e02\_h2o\_relax.vsim}]
\begin{lstlisting}[style=vsimlst]
[project]
name        = "e02_h2o_relax"
version     = "v5.14.1"
seed_base   = 42
determinism = true

[material]
formula   = "H2O"
structure = "water_molecule"   # registry alias -> SPC/E geometry
phase     = "gas"

[run]
mode         = "relax"
max_steps    = 800
converge     = true
output_level = "verbose"

[simulation]
fire_max_steps = 800
fire_dt_fs     = 0.5

[export]
write_xyz           = true
write_analysis_json = true
write_report_md     = true
output_dir          = "out/e02"

[visual]
output_type            = "terminal_chart"
render_interval        = 25
show_convergence_trace = true
show_proxy_table       = true
\end{lstlisting}
\end{tcolorbox}

\begin{termblock}
\begin{lstlisting}[style=terminallst]
  VSEPR-SIM  relax  e02_h2o_relax          step 247 / 800
  ─────────────────────────────────────────────────────
  E_total   -219.341 eV  ▼  converging
  F_max       0.0031 eV/Å ▼
  RMSD        0.0008 Å
  Energy  ╔══════════════════════════════════╗
		  ║▓▓▓▓▓░░░░░░░░░░░░░░░░░░░░░░░░░░░║  -219.3
		  ╚══════════════════════════════════╝
  CONVERGED at step 247   F_max = 0.0028 eV/Å
\end{lstlisting}
\end{termblock}

\begin{tip}
To freeze the oxygen and let only the hydrogens relax, append:
\begin{lstlisting}[style=vsimlst]
[[override.particle]]
element = "O"
fixed   = true
\end{lstlisting}
\end{tip}

%══════════════════════════════════════════════════════════════════════════════
\exsection{e03}{E03}{NaCl Crystal --- Level-0 Intent}
%══════════════════════════════════════════════════════════════════════════════

\textbf{Goal:} run a full PBC crystal relaxation with the minimum possible
authoring — let the registry fill in the details.

\textit{Key concepts:} level-0 \texttt{structure} alias · \texttt{RegistryResolver}
· PBC · \texttt{research\_report} export profile

\begin{tcolorbox}[scriptbox, scriptheader, title={e03\_nacl\_crystal.vsim}]
\begin{lstlisting}[style=vsimlst]
[project]
name        = "e03_nacl_crystal"
version     = "v5.14.1"
seed_base   = 3001
determinism = true

# Level-0: "rocksalt" alias resolved to prototype "B1_NaCl" by
# RegistryResolver. Registry fills: solver=fire, forcefield=ewald_formal,
# export_profile=research_report, observables=ionic_set.
[material]
formula   = "NaCl"
structure = "rocksalt"
cell      = "3x3x3"
phase     = "solid"

[run]
mode         = "relax"
max_steps    = 1200
converge     = true
output_level = "verbose"

[cell]
type = "orthorhombic"
lx   = 16.86    # 3 x a_NaCl (5.62 A)
ly   = 16.86
lz   = 16.86

[boundary]
x = "periodic"
y = "periodic"
z = "periodic"

[pbc]
minimum_image  = true
wrap_positions = "after_step"

[simulation]
periodic         = true
use_ewald        = true
ewald_alpha      = 0.3
ewald_rcut       = 8.0
ewald_kmax       = 5
formation_preset = "ionic"

[export]
export_profile = "research_report"
output_dir     = "out/e03"

[visual]
output_type              = "terminal_overlay_cycle"
render_interval          = 50
show_proxy_table         = true
show_convergence_trace   = true
show_rdf_plot            = true
show_steady_state_marker = true
\end{lstlisting}
\end{tcolorbox}

\noindent\textbf{Registry resolution log:}
\begin{termblock}
\begin{lstlisting}[style=terminallst]
[REGISTRY] material.prototype   <- "B1_NaCl"
[REGISTRY] material.space_group <- "Fm-3m"
[REGISTRY] run.solver           <- "fire"
[REGISTRY] run.forcefield       <- "ewald_formal"
[REGISTRY] export.profile       <- "research_report"
[REGISTRY] observe.metrics      <- "energy_map,coordination,rdf,defect"
\end{lstlisting}
\end{termblock}

\begin{keyfieldsnv}
\texttt{structure = "rocksalt"}         & Level-0: registry resolves prototype \\
\texttt{cell = "3x3x3"}                 & Supercell expansion factor \\
\texttt{export\_profile = "research\_report"} & Named profile sets 7 output flags \\
\texttt{output\_type = "terminal\_overlay\_cycle"} & Cycles RDF/energy/coordination panels \\
\end{keyfieldsnv}

%══════════════════════════════════════════════════════════════════════════════
\exsection{e04}{E04}{Silicon Diamond with RDF}
%══════════════════════════════════════════════════════════════════════════════

\textbf{Goal:} relax a Si diamond-cubic supercell, compute RDF, and visualise
the analysis overlay.

\textit{Key concepts:} covalent prototype · Tersoff FF
· \texttt{[observe]} · \texttt{[analysis.structure]}

\begin{tcolorbox}[scriptbox, scriptheader, title={e04\_silicon\_diamond.vsim}]
\begin{lstlisting}[style=vsimlst]
[project]
name        = "e04_silicon_diamond"
version     = "v5.14.1"
seed_base   = 4001
determinism = true

[material]
formula   = "Si"
prototype = "A4_Si"       # diamond cubic; registry fills Tersoff + Fd-3m
cell      = "2x2x2"
phase     = "solid"

[run]
mode      = "relax"
max_steps = 1000
converge  = true

[cell]
type = "orthorhombic"
lx   = 10.86    # 2 x a_Si (5.43 A)
ly   = 10.86
lz   = 10.86

[boundary]
x = "periodic"
y = "periodic"
z = "periodic"

[pbc]
minimum_image  = true
wrap_positions = "after_step"

[analysis.structure]
enabled              = true
compute_rdf          = true
rdf_dr               = 0.05
rdf_rmax             = 5.0
compute_coordination = true
coordination_cutoff  = 2.8

[observe]
metrics       = ["rdf", "coordination", "energy_map", "formation"]
output_format = "json"
every_n_steps = 100

[export]
write_xyz           = true
write_analysis_json = true
write_report_md     = true
output_dir          = "out/e04"

[export.visual]
write_rdf_svg     = true
visual_output_dir = "figures/e04"

[visual]
output_type            = "terminal_chart"
render_interval        = 50
show_convergence_trace = true
show_rdf_plot          = true
show_proxy_table       = true
\end{lstlisting}
\end{tcolorbox}

\begin{keyfieldsnv}
\texttt{prototype = "A4\_Si"} & Registry resolves Tersoff FF and Fd-3m space group \\
\texttt{compute\_rdf = true}  & Enables radial distribution function computation \\
\texttt{rdf\_dr = 0.05}       & Bin width in \AA{} --- smaller is finer but slower \\
\texttt{write\_rdf\_svg}      & Outputs RDF as SVG to \texttt{figures/e04/} \\
\end{keyfieldsnv}

%══════════════════════════════════════════════════════════════════════════════
\exsection{e05}{E05}{Argon Gas NVT}
%══════════════════════════════════════════════════════════════════════════════

\textbf{Goal:} simulate an Ar$_{64}$ ensemble with a thermostat and track MSD
and diffusion.

\textit{Key concepts:} \texttt{[[simulation.molecule]]} · NVT MD · MSD
· \texttt{[observe]}

\begin{tcolorbox}[scriptbox, scriptheader, title={e05\_argon\_gas\_nvt.vsim}]
\begin{lstlisting}[style=vsimlst]
[project]
name        = "e05_argon_gas"
version     = "v5.14.1"
seed_base   = 5001
determinism = true

[run]
mode          = "md"
max_steps     = 5000
dt_fs         = 2.0
temperature_K = 300.0
converge      = false

[[simulation.molecule]]
formula     = "Ar"
count       = 64
temperature = 300.0
lattice     = "random_pack"

[environment]
temperature = 300.0
pressure    = 0.0
medium      = "vacuum"

[cell]
type = "orthorhombic"
lx   = 20.0
ly   = 20.0
lz   = 20.0

[boundary]
x = "periodic"
y = "periodic"
z = "periodic"

[pbc]
minimum_image        = true
wrap_positions       = "after_step"
track_images         = true
unwrap_for_diffusion = true   # required for correct MSD

[observe]
metrics       = ["msd", "rdf", "diffusion_coefficient",
				 "temperature", "energy_map"]
output_format = "json"
every_n_steps = 50

[export]
write_xyz     = true
write_xyzf    = true
write_metrics_tsv = true
output_dir    = "out/e05"

[visual]
output_type     = "terminal_chart"
animation_mode  = "spark"
render_interval = 100
show_proxy_table = true
\end{lstlisting}
\end{tcolorbox}

\begin{keyfieldsnv}
\texttt{count = 64}                   & 64 argon atoms in the molecule block \\
\texttt{lattice = "random\_pack"}     & Initial positions from random packing \\
\texttt{unwrap\_for\_diffusion = true}& PBC images unwrapped so MSD is not truncated \\
\texttt{"diffusion\_coefficient"}     & Computes $D$ from linear MSD slope \\
\end{keyfieldsnv}

%══════════════════════════════════════════════════════════════════════════════
\exsection{e06}{E06}{TiO\textsubscript{2} Oxidation + While Loop}
%══════════════════════════════════════════════════════════════════════════════

\textbf{Goal:} run ambient oxidation MD on a TiO$_2$ rutile supercell and
continue automatically until energy variance settles.

\textit{Key concepts:} \texttt{[environment]} · \texttt{[chemistry]}
· \texttt{[variance]} · \texttt{[while]} convergence guard

\begin{tcolorbox}[scriptbox, scriptheader, title={e06\_tio2\_oxidation.vsim}]
\begin{lstlisting}[style=vsimlst]
[project]
name        = "e06_tio2_oxidation"
version     = "v5.14.1"
seed_base   = 6001
determinism = true

[material]
formula   = "TiO2"
structure = "rutile"
cell      = "2x2x2"
phase     = "solid"

[run]
mode          = "md"
max_steps     = 2000
dt_fs         = 0.5
temperature_K = 600.0
converge      = false

[cell]
type = "orthorhombic"
lx   = 9.186    # 2 x 4.593 A
ly   = 9.186
lz   = 5.918    # 2 x 2.959 A

[boundary]
x = "periodic"
y = "periodic"
z = "periodic"

[pbc]
minimum_image        = true
wrap_positions       = "after_step"
track_images         = true
unwrap_for_diffusion = true

[environment]
temperature = 600.0
pressure    = 0.0
medium      = "argon_gas"

# heat = -1: auto-derive from environment.temperature
[chemistry]
chemistry              = "oxidation"
heat                   = -1
reaction_events        = true
track_species_state    = true
min_score_threshold    = 0.25
max_reactions_per_step = 4

[observe]
metrics       = ["reaction_events", "exothermic_count", "avg_delta_E",
				 "formation", "transport", "defect"]
output_format = "json"
every_n_steps = 50

# Variance probe — guard variable used by [while]
[variance]
energy_var = "energy.total  last 100  0.05"

# Re-run 200-step blocks until energy_var <= 0.05 or 10 iterations done
[while]
name          = "settle_energy"
condition     = "variance energy_var > 0.05"
body_steps    = 200
max_iters     = 10
measure       = ["energy_var"]
iter_delay_ms = 150

[export]
write_xyz                  = true
write_xyzf                 = true
write_analysis_json        = true
write_events_json          = true
write_report_md            = true
write_dashboard_svg        = true
write_manifest_json        = true
write_pipeline_audit_jsonl = true
output_dir                 = "out/e06"

[visual]
output_type            = "terminal_chart"
animation_mode         = "spark"
render_interval        = 50
show_proxy_table       = true
show_convergence_trace = true
show_event_timeline    = true
show_rdf_plot          = true
\end{lstlisting}
\end{tcolorbox}

\noindent\textbf{While-loop execution trace:}
\begin{termblock}
\begin{lstlisting}[style=terminallst]
[while] settle_energy — iter 1/10   variance = 0.148  CONTINUE
[while] settle_energy — iter 2/10   variance = 0.091  CONTINUE
[while] settle_energy — iter 3/10   variance = 0.047  EXIT  OK
\end{lstlisting}
\end{termblock}

\begin{keyfieldsnv}
\texttt{chemistry = "oxidation"}       & Activates oxidation reaction rule family \\
\texttt{heat = -1}                     & Auto-derive reaction heat from temperature \\
\texttt{energy\_var = "... last 100 0.05"} & Probe: variance of last 100 readings, threshold 0.05 \\
\texttt{condition = "variance ... > 0.05"} & While continues while above threshold \\
\texttt{max\_iters = 10}               & Hard cap --- exits even if not converged \\
\end{keyfieldsnv}

%══════════════════════════════════════════════════════════════════════════════
\exsection{e07}{E07}{Batch Temperature Sweep}
%══════════════════════════════════════════════════════════════════════════════

\textbf{Goal:} automatically run NaCl at five temperatures and aggregate
the results into a single comparative report.

\textit{Key concepts:} \texttt{[batch]} · sweep axis · \texttt{seed\_count}
· \texttt{aggregate}

\begin{tcolorbox}[scriptbox, scriptheader, title={e07\_batch\_temperature\_sweep.vsim}]
\begin{lstlisting}[style=vsimlst]
[project]
name        = "e07_batch_sweep"
version     = "v5.14.1"
seed_base   = 7001
determinism = true

[material]
formula   = "NaCl"
prototype = "B1_NaCl"
cell      = "3x3x3"
phase     = "solid"

[run]
mode         = "md"
max_steps    = 1000
dt_fs        = 1.0
converge     = false
output_level = "minimal"

[variance]
energy_var = "energy.total  last 50  0.05"

[batch]
print_plan    = true
abort_on_fail = false

# Job 1: temperature sweep (5 temps x 1 seed = 5 sub-runs)
job         = nacl_temp_sweep
temperature = ["300", "600", "900", "1300", "2000"]
seed_count  = 1
aggregate   = true
analyze     = variance energy.total
measure     = energy_var

# Job 2: seed statistics at 900 K (1 temp x 5 seeds = 5 sub-runs)
job         = nacl_seed_stats
temperature = ["900"]
seed_count  = 5
aggregate   = true
analyze     = rmsd
measure     = energy_var

[export]
write_analysis_json = true
write_metrics_tsv   = true
write_report_md     = true
write_manifest_json = true
output_dir          = "out/e07"

[visual]
output_type      = "terminal_chart"
animation_mode   = "bar"
render_interval  = 100
show_bar_chart   = true
show_proxy_table = true
\end{lstlisting}
\end{tcolorbox}

\begin{keyfieldsnv}
\texttt{job = nacl\_temp\_sweep} & Starts a new named job entry \\
\texttt{temperature = ["300","600",...]} & Sweep axis --- one sub-run per value \\
\texttt{seed\_count = 5}         & Repeats each parameter point with different seeds \\
\texttt{aggregate = true}        & Merges sub-run results into a single report \\
\end{keyfieldsnv}

%══════════════════════════════════════════════════════════════════════════════
\exsection{e08}{E08}{GL Visualizer Controls}
%══════════════════════════════════════════════════════════════════════════════

\textbf{Goal:} demonstrate all major interactive GL visualizer options.

\textit{Key concepts:} \texttt{gl\_interactive} · spin · orbit · overlay cycle

\begin{tcolorbox}[scriptbox, scriptheader, title={e08\_gl\_visualizer.vsim}]
\begin{lstlisting}[style=vsimlst]
[project]
name    = "e08_gl_visualizer"
version = "v5.14.1"

[material]
formula   = "Cu"
structure = "fcc"
cell      = "3x3x3"
phase     = "solid"

[run]
mode      = "relax"
max_steps = 500
converge  = true

[export]
write_xyz  = true
output_dir = "out/e08"

[visual]
output_type       = "gl_interactive"    # full ImGui viewer
gl_window_width   = 1280
gl_window_height  = 800
gl_auto_orbit     = false
gl_spin           = true
gl_spin_axis      = "y"
gl_spin_deg_per_s = 20.0
gl_show_axes      = true
render_interval   = 10
\end{lstlisting}
\end{tcolorbox}

\noindent\textbf{\texttt{output\_type} comparison:}

\begin{center}
\footnotesize
\begin{tabular}{@{}llll@{}}
\toprule
\texttt{output\_type} & Window? & Interactive? & Best for \\
\midrule
\texttt{"none"}              & No  & No      & Headless batch, CI \\
\texttt{"terminal\_chart"}   & No  & ASCII   & Live monitoring, SSH \\
\texttt{"gl\_interactive"}   & Yes & Yes     & 3D inspection, ImGui \\
\texttt{"gl\_live\_60fps"}   & Yes & No      & Smooth demos \\
\texttt{"gl\_crystal\_grid"} & Yes & No      & Crystal structures \\
\texttt{"gl\_overlay\_cycle"}& Yes & No      & Cycling analysis panels \\
\bottomrule
\end{tabular}
\end{center}

\begin{tip}
\texttt{gl\_spin} and \texttt{gl\_auto\_orbit} are mutually exclusive.
Setting \texttt{gl\_spin = true} disables \texttt{gl\_auto\_orbit} for that run.
Use \texttt{gl\_spin\_deg\_per\_s = 0.0} for a static frozen view.
\end{tip}

%══════════════════════════════════════════════════════════════════════════════
\exsection{e09}{E09}{IKK Identity Vector}
%══════════════════════════════════════════════════════════════════════════════

\textbf{Goal:} enable the IKK identity-vector analysis to output a
\texttt{.identity.json} file tracking $\bar{I}_f$ across frames.

\textit{Key concepts:} \texttt{[analysis.ivec]} · I-vector axes
· \texttt{.identity.json} · WO-75B Phase 1

\begin{tcolorbox}[scriptbox, scriptheader, title={e09\_ikk\_identity\_vector.vsim}]
\begin{lstlisting}[style=vsimlst]
[project]
name    = "nacl_ikk"
version = "v5.14.1"

[material]
formula   = "NaCl"
structure = "rocksalt"
cell      = "2x2x2"
phase     = "solid"

[run]
mode      = "md"
max_steps = 500
dt_fs     = 1.0
converge  = false

[cell]
type = "orthorhombic"
lx   = 11.24
ly   = 11.24
lz   = 11.24

[boundary]
x = "periodic"
y = "periodic"
z = "periodic"

[pbc]
minimum_image  = true
wrap_positions = "after_step"

# IKK identity-vector analysis (WO-75B Phase 1)
# Phase 1 proxy mapping (one aggregate per frame):
#   x = 1 - dataloss       (existence)
#   y = 1 - hidden_channel (EM)
#   z = recoverable_info   (spatial)
#   t = 1 - projection_loss(temporal)
#   w = 1-identity_residual(internal)
[analysis.ivec]
enabled       = true
write_json    = true
include_delta = true
include_var   = false
output_dir    = "out/e09"

[analysis.ikk_end_tag]
enabled           = true
emit_markdown     = true
d_pass_threshold  = 0.65
d_warn_threshold  = 0.40
scale_regime      = "S3"
section_reference = "IV.9"

[export]
write_xyz       = true
write_report_md = true
output_dir      = "out/e09"

[visual]
output_type      = "terminal_chart"
render_interval  = 50
show_proxy_table = true
\end{lstlisting}
\end{tcolorbox}

\noindent\textbf{I-vector axis mapping (Phase 1):}

\begin{center}
\footnotesize
\begin{tabular}{@{}cll@{}}
\toprule
Axis & Sidecar field & Interpretation \\
\midrule
$x$ & $1 - \texttt{dataloss}$          & Existence: identity preservation fraction \\
$y$ & $1 - \texttt{hidden\_channel}$   & EM: complement of hidden-interaction proxy \\
$z$ & $\texttt{recoverable\_info}$      & Spatial: recoverable fraction \\
$t$ & $1 - \texttt{projection\_loss}$  & Temporal: complement of scale-crossing loss \\
$w$ & $1 - \texttt{identity\_residual}$& Internal: complement of prototype deviation \\
\bottomrule
\end{tabular}
\end{center}

Axis order $(x,y,z,t,w)$ is \textbf{fixed and non-permutable}.

\begin{doctrine}
\texttt{[analysis.ivec]} is derived data only.
It reads from \texttt{IdentitySidecarSeries} and never writes back to
\texttt{.xyz} or \texttt{.xyzFull} truth-state files.
\end{doctrine}

%══════════════════════════════════════════════════════════════════════════════
\exsection{e10}{E10}{MCF-CAI Object State Grid}
%══════════════════════════════════════════════════════════════════════════════

\textbf{Goal:} populate the $3\times3$ MCF-CAI object state grid for a carbon
diamond bead.

\textit{Key concepts:} \texttt{[object.<layer>.<basis>]}
· Information-column doctrine · WO-76

\begin{tcolorbox}[scriptbox, scriptheader, title={e10\_mcf\_cai.vsim}]
\begin{lstlisting}[style=vsimlst]
# Grid: rows = macro|chemical|fundamental
#       cols = carrier|action|information (information is sidecar-only)

[project]
name    = "e10_mcf_cai"
version = "v5.14.1"

[material]
formula   = "C"
prototype = "A4_diamond"
cell      = "1x1x1"
phase     = "solid"

[run]
mode      = "relax"
max_steps = 400
converge  = true

[export]
write_xyz  = true
output_dir = "out/e10"

[visual]
output_type     = "terminal_chart"
render_interval = 50

# ── MACRO row ─────────────────────────────────────────────────
[object.macro.carrier]
centre_x    = 0.0
phase_label = "diamond_cubic"
grain_count = 1

[object.macro.action]
stress_proxy    = 0.0
fracture_prob   = 0.0

[object.macro.information]         # sidecar only
formation_age_fs = 0.0
defect_memory    = 0.0

# ── CHEMICAL row ──────────────────────────────────────────────
[object.chemical.carrier]
atomic_number       = 6            # Carbon
mass                = 12.011
coordination_number = 4            # sp3

[object.chemical.action]
reaction_active = false
oxidation_state = 0

[object.chemical.information]      # sidecar only
bond_entropy_loss = 0.0
d_chem            = 0.0
formation_route   = "tetrahedral_sp3"

# ── FUNDAMENTAL row ────────────────────────────────────────────
[object.fundamental.carrier]
net_charge    = 0.0
spin_proxy    = 0.0
isotope_label = "C12"

[object.fundamental.action]
em_coupling  = 0.0
strong_proxy = 1.0

[object.fundamental.information]   # sidecar only
psi_hid         = 0.0
projection_loss = 0.0
entropy_trace   = 0.0
\end{lstlisting}
\end{tcolorbox}

\noindent\textbf{MCF-CAI grid:}

\begin{center}
\footnotesize
\begin{tabular}{@{}l|l|l|l@{}}
\toprule
 & \textbf{carrier} & \textbf{action} & \textbf{information (sidecar)} \\
\midrule
\textbf{macro}       & centre, phase\_label, grain\_count & stress, heat\_flux, fracture & formation\_age\_fs, defect\_memory \\
\textbf{chemical}    & Z, mass, CN, bond\_order          & reaction, oxidation, ionisation & bond\_entropy, d\_chem, route \\
\textbf{fundamental} & net\_charge, spin, isotope        & em, strong, weak, decay & psi\_hid, projection\_loss, entropy \\
\bottomrule
\end{tabular}
\end{center}

\begin{doctrine}
\texttt{information}-column cells are sidecar/audit data only.
Their fields must \textbf{never} be passed to the dynamics force kernel and
must \textbf{never} be written back to truth-state files (\texttt{.xyz} /
\texttt{.xyzFull}).
\end{doctrine}

%══════════════════════════════════════════════════════════════════════════════
\exsection{e11}{E11}{Full Research Pipeline}
%══════════════════════════════════════════════════════════════════════════════

\textbf{Goal:} combine every major section into one production-grade research
run on FeO w\"{u}stite.

\textit{Key concepts:} every section working together
· \texttt{publication} export profile · IKK end-tag · MCF-CAI

\begin{tcolorbox}[scriptbox, scriptheader, title={e11\_full\_pipeline.vsim}]
\begin{lstlisting}[style=vsimlst]
[project]
name        = "e11_feo_full_pipeline"
version     = "v5.14.1"
seed_base   = 11001
determinism = true

[material]
formula   = "FeO"
structure = "wustite"
cell      = "2x2x2"
phase     = "solid"

[run]
mode          = "md"
max_steps     = 3000
dt_fs         = 0.5
temperature_K = 800.0
converge      = false
output_level  = "verbose"

[cell]
type = "orthorhombic"
lx   = 8.60
ly   = 8.60
lz   = 8.60

[boundary]
x = "periodic"
y = "periodic"
z = "periodic"

[pbc]
minimum_image        = true
wrap_positions       = "after_step"
track_images         = true
unwrap_for_diffusion = true

[simulation]
periodic         = true
use_ewald        = true
ewald_alpha      = 0.3
ewald_rcut       = 7.0
ewald_kmax       = 5
formation_preset = "ionic"

[environment]
temperature = 800.0
pressure    = 0.0
medium      = "argon_gas"

[chemistry]
chemistry              = "oxidation"
heat                   = -1
reaction_events        = true
track_species_state    = true
min_score_threshold    = 0.20
max_reactions_per_step = 6

[observe]
metrics       = ["rdf", "coordination", "msd", "energy_map",
				 "reaction_events", "defect", "transport"]
output_format = "json"
every_n_steps = 100

[analysis.structure]
enabled              = true
compute_rdf          = true
rdf_dr               = 0.05
rdf_rmax             = 6.0
compute_coordination = true
coordination_cutoff  = 2.6

[analysis.ikk_end_tag]
enabled           = true
emit_markdown     = true
emit_latex        = true
d_pass_threshold  = 0.65
d_warn_threshold  = 0.40
scale_regime      = "S3"
section_reference = "IV.11"

[analysis.ivec]
enabled       = true
write_json    = true
include_delta = true
output_dir    = "out/e11"

[object.macro.carrier]
phase_label = "wustite_B1"
grain_count = 8

[object.chemical.carrier]
atomic_number       = 26
mass                = 55.845
coordination_number = 6

[object.chemical.action]
reaction_active = true
oxidation_state = 2

[object.fundamental.carrier]
net_charge    = 2.0
spin_proxy    = 2.0
isotope_label = "Fe56"

[object.fundamental.information]
psi_hid         = 0.0
projection_loss = 0.0
entropy_trace   = 0.0

[variance]
energy_var = "energy.total  last 150  0.04"

[while]
name          = "settle_energy"
condition     = "variance energy_var > 0.04"
body_steps    = 300
max_iters     = 8
measure       = ["energy_var"]
iter_delay_ms = 200

[export]
export_profile             = "publication"
write_pipeline_audit_jsonl = true
output_dir                 = "out/e11"

[export.visual]
write_svg_figures         = true
write_rdf_svg             = true
write_energy_trace_svg    = true
write_packing_heatmap_svg = true
visual_output_dir         = "figures/e11"

[visual]
output_type              = "terminal_overlay_cycle"
animation_mode           = "spark"
render_interval          = 50
show_proxy_table         = true
show_convergence_trace   = true
show_event_timeline      = true
show_rdf_plot            = true
show_steady_state_marker = true
show_audit_table         = true

[report]
title                  = "FeO Wustite - Full Research Pipeline"
include_material_cards = true
include_metric_tables  = true
\end{lstlisting}
\end{tcolorbox}

%══════════════════════════════════════════════════════════════════════════════
\newpage
\section*{QR — Quick-Reference Card}
\addcontentsline{toc}{section}{QR — Quick-Reference Card}
%══════════════════════════════════════════════════════════════════════════════

\subsection*{Mandatory fields}
\begin{tcolorbox}[scriptbox]
\begin{lstlisting}[style=vsimlst]
[project]
name    = "<name>"      # required
version = "v5.14.1"     # recommended

[material]
formula = "<formula>"   # required; everything else can be registry-filled

[run]
mode    = "<relax|md>"  # required
\end{lstlisting}
\end{tcolorbox}

\subsection*{Common \texttt{[run]} patterns}

\begin{center}
\footnotesize
\begin{tabular}{@{}ll@{}}
\toprule
Intent & Fields \\
\midrule
Energy minimisation  & \texttt{mode="relax"}, \texttt{max\_steps=1000}, \texttt{converge=true} \\
NVT molecular dynamics & \texttt{mode="md"}, \texttt{dt\_fs=1.0}, \texttt{temperature\_K=300.0}, \texttt{converge=false} \\
Static display only  & \texttt{mode="md"}, \texttt{max\_steps=1}, \texttt{converge=false} \\
\bottomrule
\end{tabular}
\end{center}

\subsection*{Export profiles}

\begin{center}
\footnotesize
\begin{tabular}{@{}ll@{}}
\toprule
Profile & Files emitted \\
\midrule
\texttt{minimal}         & \texttt{.xyz} \\
\texttt{standard}        & \texttt{.xyz} + \texttt{.analysis.json} + \texttt{.metrics.tsv} + \texttt{.report.md} \\
\texttt{research\_report}& standard + events, manifest, dashboard SVG \\
\texttt{publication}     & research\_report + symbolic trace + pipeline audit JSONL \\
\bottomrule
\end{tabular}
\end{center}

\subsection*{PBC quick setup}
\begin{tcolorbox}[scriptbox]
\begin{lstlisting}[style=vsimlst]
[cell]
type = "orthorhombic"
lx = 10.0  ly = 10.0  lz = 10.0

[boundary]
x = "periodic"  y = "periodic"  z = "periodic"

[pbc]
minimum_image  = true
wrap_positions = "after_step"
\end{lstlisting}
\end{tcolorbox}

\subsection*{Analysis section summary}

\begin{center}
\footnotesize
\begin{tabular}{@{}lll@{}}
\toprule
Section & Purpose & WO \\
\midrule
\texttt{[analysis.structure]}       & RDF, coordination number          & --- \\
\texttt{[analysis.sampling]}        & Scalar sampling, MSD              & --- \\
\texttt{[analysis.scale\_sampling]} & Field projection, RVE, emergence  & WO-61D \\
\texttt{[analysis.inference]}       & Property inference                & WO-61B \\
\texttt{[analysis.ikk\_end\_tag]}   & IKK report end-tag badge          & WO-75A \\
\texttt{[analysis.ivec]}            & I-vector series $\to$ \texttt{.identity.json} & WO-75B \\
\texttt{[object.<layer>.<basis>]}   & MCF-CAI $3\times3$ object state   & WO-76  \\
\texttt{[object.surface]}           & Surface analysis master block     & Surface \\
\texttt{[[object.surface]]}         & Per-surface probe definition      & Surface \\
\bottomrule
\end{tabular}
\end{center}

%==============================================================================
\newpage
\section*{Surface Analysis Examples (E12--E27)}
%==============================================================================

\begin{tcolorbox}[
  enhanced, breakable,
  colback=vsimdoctrinebg, colframe=vsimdoctrinetop,
  arc=3pt, boxrule=1pt,
  fonttitle=\sffamily\bfseries\small\color{white},
  title={Doctrine: Analysis-Only Probes},
]
\texttt{[[object.surface]]} sections are \textbf{measurement-only}.
They never modify particle positions, velocities, or forces, and they
never touch the canonical \texttt{.xyz} truth-state.
All metrics land in a \texttt{.surface.json} sidecar keyed by \texttt{output\_tag}.
\end{tcolorbox}

\medskip
\begin{center}\footnotesize
\begin{tabular}{@{}cll@{}}
\toprule
Symbol & Field & Meaning \\
\midrule
$\Phi$   & \texttt{compute\_flux}           & Net particle crossings / (area\,$\cdot$\,step) \\
$Q$      & \texttt{compute\_energy\_flux}   & Kinetic energy flux / (area\,$\cdot$\,step) \\
$P$      & \texttt{compute\_momentum\_flux} & Normal momentum transfer / (area\,$\cdot$\,step) \\
$J_s$    & \texttt{compute\_species\_flux}  & Per-species crossing table \\
$J_m$    & \texttt{compute\_mass\_flux}     & Mass-weighted crossing flux \\
$J_{iv}$ & \texttt{compute\_ivec\_flux}     & IVec crossing signature (needs \texttt{[analysis.ivec]}) \\
\bottomrule
\end{tabular}
\end{center}

\bigskip

%------------------------------------------------------------------------------
\subsection*{E12 --- Minimal Rectangular Surface Probe}
%------------------------------------------------------------------------------

\textbf{Goal:} simplest possible \texttt{[[object.surface]]} usage ---
a rectangle counts particle crossings at the XZ mid-plane, no flux metric.\\
\textbf{Key concepts:} \texttt{[object.surface]} master block \textbullet{}
\texttt{[[object.surface]]} per-probe entry \textbullet{}
\texttt{log\_crossings} \textbullet{} analysis-only doctrine

\begin{tcolorbox}[scriptbox, scriptheader,
  title={scripts/e12\_surface\_probe\_minimal.vsim}]
\begin{lstlisting}[style=vsimlst]
[object.surface]
enabled    = true
write_json = true
output_dir = "out/e12/surfaces"

[[object.surface]]
name          = "midplane_XZ"
geometry      = "rectangle"
center_x      = 0.0
center_y      = 0.0
center_z      = 0.0
normal_x      = 0.0
normal_y      = 1.0
normal_z      = 0.0
width         = 20.0
height        = 20.0
log_crossings = true
compute_flux  = false
output_tag    = "surf_midXZ"
\end{lstlisting}
\end{tcolorbox}

%------------------------------------------------------------------------------
\subsection*{E13 --- Particle Number Flux ($\Phi$)}
%------------------------------------------------------------------------------

\textbf{Goal:} compute net signed particle flux $\Phi$ through a rectangular
channel probe in an Ar gas.\\
\textbf{Key concepts:} \texttt{compute\_flux = true} \textbullet{} $\Phi$ =
net crossings / (area\,$\cdot$\,step) \textbullet{} \texttt{.surface.json} sidecar

\begin{tcolorbox}[scriptbox, scriptheader,
  title={scripts/e13\_surface\_particle\_flux.vsim}]
\begin{lstlisting}[style=vsimlst]
[object.surface]
enabled    = true
write_json = true
output_dir = "out/e13/surfaces"

[[object.surface]]
name          = "channel_probe"
geometry      = "rectangle"
center_x      = 0.0
center_y      = 0.0
center_z      = 0.0
normal_x      = 1.0
normal_y      = 0.0
normal_z      = 0.0
width         = 30.0
height        = 30.0
log_crossings = true
compute_flux  = true
output_tag    = "phi_channel"
\end{lstlisting}
\end{tcolorbox}

%------------------------------------------------------------------------------
\subsection*{E14 --- Time-Averaged Flux Statistics}
%------------------------------------------------------------------------------

\textbf{Goal:} compute rolling time-averaged flux over a long MD run to identify steady-state vs fluctuating flow.\\
\textbf{Key concepts:} \texttt{compute\_flux} per-step data \textbullet{}
temporal statistics \textbullet{} quasi-steady-state detection

\begin{tcolorbox}[scriptbox, scriptheader,
  title={scripts/e14\_surface\_time\_averaged\_flux.vsim (excerpt)}]
\begin{lstlisting}[style=vsimlst]
[object.surface]
enabled    = true
write_json = true
output_dir = "out/e14/surfaces"

[[object.surface]]
name          = "flux_temporal_probe"
geometry      = "rectangle"
normal_y      = 1.0
width         = 25.0
height        = 25.0
log_crossings = true
compute_flux  = true
output_tag    = "phi_temporal"
\end{lstlisting}
\end{tcolorbox}

\begin{tcolorbox}[enhanced, colback=vsimcodebg, colframe=vsimrule,
  arc=3pt, boxrule=0.6pt, fonttitle=\sffamily\small\bfseries\color{vsimkw},
  title={Tip}]
\footnotesize
Extract the \texttt{flux} array from \texttt{.surface.json} and compute
rolling mean/variance over a sliding window (e.g.\ 500 steps) to detect
quasi-steady-state.
\end{tcolorbox}

%------------------------------------------------------------------------------
\subsection*{E15 --- Species-Filtered Flux}
%------------------------------------------------------------------------------

\textbf{Goal:} two co-located probes record Na$^+$ and Cl$^-$ flux separately.\\
\textbf{Key concepts:} \texttt{species\_filter} \textbullet{}
\texttt{compute\_mass\_flux} \textbullet{} dual-probe comparative output

\begin{tcolorbox}[scriptbox, scriptheader,
  title={scripts/e15\_surface\_species\_filtered\_flux.vsim}]
\begin{lstlisting}[style=vsimlst]
[object.surface]
enabled    = true
write_json = true
output_dir = "out/e15/surfaces"

[[object.surface]]          # Na-only probe
name              = "na_membrane"
geometry          = "rectangle"
normal_y          = 1.0
width             = 25.0
height            = 20.0
log_crossings     = true
compute_flux      = true
compute_mass_flux = true
species_filter    = ["Na"]
output_tag        = "phi_Na"

[[object.surface]]          # Cl-only probe
name              = "cl_membrane"
geometry          = "rectangle"
normal_y          = 1.0
width             = 25.0
height            = 20.0
log_crossings     = true
compute_flux      = true
compute_mass_flux = true
species_filter    = ["Cl"]
output_tag        = "phi_Cl"
\end{lstlisting}
\end{tcolorbox}

\begin{tcolorbox}[enhanced, colback=vsimcodebg, colframe=vsimrule,
  arc=3pt, boxrule=0.6pt, fonttitle=\sffamily\small\bfseries\color{vsimkw},
  title={Tip}]
\footnotesize
\texttt{species\_filter} accepts a TOML inline array
\texttt{["Na",\,"Cl"]} or a single bare string \texttt{"Na"}.
\end{tcolorbox}

%------------------------------------------------------------------------------
\subsection*{E16 --- Kinetic Energy Flux ($Q$)}
%------------------------------------------------------------------------------

\textbf{Goal:} measure $Q$ through a disk probe in a hot Ar gas.\\
\textbf{Key concepts:} \texttt{geometry = "disk"} \textbullet{}
\texttt{radius} \textbullet{} \texttt{compute\_energy\_flux} \textbullet{} thermal plane

\begin{tcolorbox}[scriptbox, scriptheader,
  title={scripts/e16\_surface\_energy\_flux.vsim}]
\begin{lstlisting}[style=vsimlst]
[object.surface]
enabled    = true
write_json = true
output_dir = "out/e16/surfaces"

[[object.surface]]
name                = "thermal_disk_z"
geometry            = "disk"
center_z            = 5.0
normal_z            = 1.0
radius              = 12.0
log_crossings       = true
compute_energy_flux = true
output_tag          = "Q_disk_z"
\end{lstlisting}
\end{tcolorbox}

\begin{center}\footnotesize
\begin{tabular}{@{}ll@{}}
\toprule
Field & Disk-specific notes \\
\midrule
\texttt{geometry = "disk"} & Circular probe; \texttt{radius} required \\
\texttt{width} / \texttt{height} & Ignored for disk geometry \\
\texttt{radius} & Half-diameter in \AA \\
\bottomrule
\end{tabular}
\end{center}

%------------------------------------------------------------------------------
\subsection*{E17 --- Multi-Metric Single Probe}
%------------------------------------------------------------------------------

\textbf{Goal:} a single probe computes mass and energy flux simultaneously on one plane.\\
\textbf{Key concepts:} \texttt{compute\_mass\_flux + compute\_energy\_flux}
\textbullet{} combined metrics in one sidecar \textbullet{} analysis locality

\begin{tcolorbox}[scriptbox, scriptheader,
  title={scripts/e17\_surface\_multi\_metric\_probe.vsim (excerpt)}]
\begin{lstlisting}[style=vsimlst]
[object.surface]
enabled    = true
write_json = true
output_dir = "out/e17/surfaces"

[[object.surface]]
name                = "multi_metric_probe"
geometry            = "rectangle"
normal_x            = 1.0
width               = 30.0
height              = 30.0
log_crossings       = true
compute_flux        = true
compute_mass_flux   = true
compute_energy_flux = true
output_tag          = "multi_metric"
\end{lstlisting}
\end{tcolorbox}

\begin{tcolorbox}[enhanced, colback=vsimcodebg, colframe=vsimrule,
  arc=3pt, boxrule=0.6pt, fonttitle=\sffamily\small\bfseries\color{vsimkw},
  title={Tip}]
\footnotesize
A single \texttt{output\_tag} holds all requested metrics; the sidecar entry
contains \texttt{"flux"}, \texttt{"mass\_flux"}, and \texttt{"energy\_flux"}
keys.
\end{tcolorbox}

%------------------------------------------------------------------------------
\subsection*{E18 --- Momentum Flux / Pressure Proxy ($P$)}
%------------------------------------------------------------------------------

\textbf{Goal:} estimate the kinetic pressure gradient across an Ar slab
using top and bottom probes at $\pm15$\,\AA.\\
\textbf{Key concepts:} \texttt{compute\_momentum\_flux} \textbullet{}
signed normal-velocity weighting \textbullet{} $\Delta P =
P_\mathrm{top} - P_\mathrm{bottom}$

\begin{tcolorbox}[scriptbox, scriptheader,
  title={scripts/e18\_surface\_momentum\_flux.vsim}]
\begin{lstlisting}[style=vsimlst]
[[object.surface]]
name                  = "slab_top"
geometry              = "rectangle"
center_y              = 15.0
normal_y              = 1.0
width                 = 30.0
height                = 30.0
log_crossings         = true
compute_flux          = true
compute_momentum_flux = true
output_tag            = "P_top"

[[object.surface]]
name                  = "slab_bottom"
geometry              = "rectangle"
center_y              = -15.0
normal_y              = 1.0
width                 = 30.0
height                = 30.0
log_crossings         = true
compute_flux          = true
compute_momentum_flux = true
output_tag            = "P_bottom"
\end{lstlisting}
\end{tcolorbox}

%------------------------------------------------------------------------------
\subsection*{E19 --- Per-Species Flux Probes}
%------------------------------------------------------------------------------

\textbf{Goal:} three co-located probes separately record Na, Cl, and Ar
crossings in a mixed simulation.\\
\textbf{Key concepts:} \texttt{compute\_species\_flux} \textbullet{}
per-species sidecar subtables \textbullet{} \texttt{output\_tag} keying

\begin{tcolorbox}[scriptbox, scriptheader,
  title={scripts/e19\_surface\_species\_flux\_multi.vsim (excerpt)}]
\begin{lstlisting}[style=vsimlst]
[[object.surface]]
name                 = "probe_Na"
geometry             = "rectangle"
normal_z             = 1.0
width                = 20.0
height               = 20.0
log_crossings        = true
compute_species_flux = true
species_filter       = ["Na"]
output_tag           = "J_Na"

[[object.surface]]
name                 = "probe_Cl"
geometry             = "rectangle"
normal_z             = 1.0
width                = 20.0
height               = 20.0
log_crossings        = true
compute_species_flux = true
species_filter       = ["Cl"]
output_tag           = "J_Cl"
\end{lstlisting}
\end{tcolorbox}

%------------------------------------------------------------------------------
\subsection*{E20 --- Identity-Vector (IVec) Flux}
%------------------------------------------------------------------------------

\textbf{Goal:} couple surface probes to the IVec module (WO-75B) at a Si
crystal boundary disk.\\
\textbf{Key concepts:} \texttt{compute\_ivec\_flux} \textbullet{}
\texttt{[analysis.ivec]} co-activation \textbullet{} WO-75B

\begin{tcolorbox}[scriptbox, scriptheader,
  title={scripts/e20\_surface\_ivec\_flux.vsim}]
\begin{lstlisting}[style=vsimlst]
[analysis.ivec]
enabled         = true
write_json      = true
convergence_eps = 1.0e-4

[object.surface]
enabled    = true
write_json = true
output_dir = "out/e20/surfaces"

[[object.surface]]
name              = "ivec_boundary_disk"
geometry          = "disk"
center_z          = 0.0
normal_z          = 1.0
radius            = 8.0
log_crossings     = true
compute_ivec_flux = true
output_tag        = "ivec_surf_z"
\end{lstlisting}
\end{tcolorbox}

\begin{tcolorbox}[enhanced, colback=vsimdoctrinebg, colframe=vsimdoctrinetop,
  arc=3pt, boxrule=0.8pt, left=6pt, right=6pt, top=4pt, bottom=4pt]
\footnotesize
\textbf{Coupling rule:} \texttt{compute\_ivec\_flux = true} is silently
ignored unless \texttt{[analysis.ivec] enabled = true} is also present
in the same script.
\end{tcolorbox}

%------------------------------------------------------------------------------
\subsection*{E21 --- Spherical Surface Geometry}
%------------------------------------------------------------------------------

\textbf{Goal:} measure flux across a closed spherical shell to model a droplet
or bubble interface.\\
\textbf{Key concepts:} \texttt{geometry = "sphere"} \textbullet{}
closed 3D shell with \texttt{radius} \textbullet{} isotropic flux \textbullet{}
droplet/bubble diagnostics

\begin{tcolorbox}[scriptbox, scriptheader,
  title={scripts/e21\_surface\_sphere\_geometry.vsim (excerpt)}]
\begin{lstlisting}[style=vsimlst]
[object.surface]
enabled    = true
write_json = true
output_dir = "out/e21/surfaces"

[[object.surface]]
name                = "droplet_shell"
geometry            = "sphere"
center_x            = 0.0
center_y            = 0.0
center_z            = 0.0
radius              = 10.0
log_crossings       = true
compute_flux        = true
compute_energy_flux = true
output_tag          = "sphere_evap"
\end{lstlisting}
\end{tcolorbox}

\begin{tcolorbox}[enhanced, colback=vsimcodebg, colframe=vsimrule,
  arc=3pt, boxrule=0.6pt, fonttitle=\sffamily\small\bfseries\color{vsimkw},
  title={Sphere note}]
\footnotesize
\texttt{geometry = "sphere"} activates a closed shell; positive flux = outward
crossing (evaporation). Area = $4\pi r^2$.
\end{tcolorbox}

%------------------------------------------------------------------------------
\subsection*{E22 --- Multi-Probe Flux Pipeline}
%------------------------------------------------------------------------------

\textbf{Goal:} combine $\Phi$, $Q$, and $P$ probes in one script;
build a three-metric sidecar for a 2\,048-atom Ar gas.\\
\textbf{Key concepts:} three \texttt{[[object.surface]]} blocks \textbullet{}
\texttt{write\_report\_md} companion \textbullet{} combined \texttt{.surface.json}

\begin{tcolorbox}[scriptbox, scriptheader,
  title={scripts/e22\_surface\_multi\_probe\_pipeline.vsim (surface block)}]
\begin{lstlisting}[style=vsimlst]
[object.surface]
enabled    = true
write_json = true
output_dir = "out/e22/surfaces"

# Phi probe at Z = -10
[[object.surface]]
name          = "phi_probe_bottom"
geometry      = "rectangle"
center_z      = -10.0
normal_z      = 1.0
width         = 40.0
height        = 40.0
log_crossings = true
compute_flux  = true
output_tag    = "phi_z_neg10"

# Q probe (disk) at Z = 0
[[object.surface]]
name                = "Q_probe_mid"
geometry            = "disk"
center_z            = 0.0
normal_z            = 1.0
radius              = 15.0
log_crossings       = true
compute_energy_flux = true
output_tag          = "Q_disk_z0"

# P probe at Z = +10
[[object.surface]]
name                  = "P_probe_top"
geometry              = "rectangle"
center_z              = 10.0
normal_z              = 1.0
width                 = 40.0
height                = 40.0
log_crossings         = true
compute_momentum_flux = true
output_tag            = "P_z_pos10"
\end{lstlisting}
\end{tcolorbox}

%------------------------------------------------------------------------------
\subsection*{E23 --- Batch Parameter Sweep with Surface Probes}
%------------------------------------------------------------------------------

\textbf{Goal:} integrate surface analysis with \texttt{[batch]} parameter
expansion; sweep temperature \{300, 400, 500, 600\} K and record flux.\\
\textbf{Key concepts:} \texttt{[batch]} matrix \textbullet{} per-job
\texttt{.surface.json} \textbullet{} aggregated flux-vs-temperature post-analysis

\begin{tcolorbox}[scriptbox, scriptheader,
  title={scripts/e23\_surface\_batch\_sweep.vsim (excerpt)}]
\begin{lstlisting}[style=vsimlst]
[batch]
enabled = true
mode    = "matrix"

[[batch.parameter]]
section = "simulation.molecule"
key     = "temperature"
values  = [300.0, 400.0, 500.0, 600.0]

[object.surface]
enabled    = true
write_json = true
output_dir = "out/e23/surfaces"

[[object.surface]]
name          = "batch_flux_probe"
geometry      = "rectangle"
normal_z      = 1.0
width         = 30.0
height        = 30.0
log_crossings = true
compute_flux  = true
output_tag    = "phi_batch"
\end{lstlisting}
\end{tcolorbox}

%------------------------------------------------------------------------------
\subsection*{E24 --- While-Loop Convergence on Flux Stationarity}
%------------------------------------------------------------------------------

\textbf{Goal:} couple \texttt{[while]} convergence guard to surface flux
variance; continue until quasi-steady-state flow.\\
\textbf{Key concepts:} \texttt{[while]} criteria on flux variance
\textbullet{} iterative simulation gated on analysis \textbullet{}
stationarity detection

\begin{tcolorbox}[scriptbox, scriptheader,
  title={scripts/e24\_surface\_while\_loop\_convergence.vsim (excerpt)}]
\begin{lstlisting}[style=vsimlst]
[variance]
name         = "flux_variance"
target       = 0.0001
tolerance    = 0.00005
check_every  = 200
max_restarts = 5

[while]
criteria    = "flux_variance"
max_cycles  = 10
cycle_steps = 1000

[object.surface]
enabled    = true
write_json = true
output_dir = "out/e24/surfaces"

[[object.surface]]
name          = "flux_convergence_probe"
geometry      = "rectangle"
normal_y      = 1.0
width         = 25.0
height        = 25.0
log_crossings = true
compute_flux  = true
output_tag    = "phi_while"
\end{lstlisting}
\end{tcolorbox}

%------------------------------------------------------------------------------
\subsection*{E25 --- Surface Flux + RDF Correlation Analysis}
%------------------------------------------------------------------------------

\textbf{Goal:} simultaneously activate \texttt{[analysis.structure]} (RDF)
and surface probes to correlate flux with local coordination.\\
\textbf{Key concepts:} \texttt{compute\_rdf = true} $\to$ \texttt{.rdf.json}
\textbullet{} \texttt{compute\_flux = true} $\to$ \texttt{.surface.json}
\textbullet{} cross-module coupling \textbullet{} liquid--solid interface

\begin{tcolorbox}[scriptbox, scriptheader,
  title={scripts/e25\_surface\_rdf\_correlation.vsim (excerpt)}]
\begin{lstlisting}[style=vsimlst]
[analysis.structure]
enabled     = true
compute_rdf = true
rdf_rmax    = 10.0
rdf_nbins   = 200

[object.surface]
enabled    = true
write_json = true
output_dir = "out/e25/surfaces"

[[object.surface]]
name          = "interface_probe"
geometry      = "rectangle"
normal_z      = 1.0
width         = 20.0
height        = 20.0
log_crossings = true
compute_flux  = true
output_tag    = "phi_rdf"
\end{lstlisting}
\end{tcolorbox}

%------------------------------------------------------------------------------
\subsection*{E26 --- Surface Flux Verification with [verify.structure]}
%------------------------------------------------------------------------------

\textbf{Goal:} combine surface analysis with the empirical verification layer
(WO-62A) to check flux consistency with mass conservation.\\
\textbf{Key concepts:} \texttt{[verify.structure]} + surface probes
\textbullet{} analysis $\to$ verification pipeline \textbullet{} co-activation

\begin{tcolorbox}[scriptbox, scriptheader,
  title={scripts/e26\_surface\_flux\_verification.vsim (excerpt)}]
\begin{lstlisting}[style=vsimlst]
[verify]
enabled = true

[verify.structure]
check_rdf                = true
rdf_first_peak_min       = 3.0
rdf_first_peak_max       = 4.0
rdf_coordination_min     = 10.0
rdf_coordination_max     = 14.0
rdf_coordination_species = "Ne"

[object.surface]
enabled    = true
write_json = true
output_dir = "out/e26/surfaces"

[[object.surface]]
name          = "verify_flux_probe"
geometry      = "rectangle"
normal_x      = 1.0
width         = 25.0
height        = 25.0
log_crossings = true
compute_flux  = true
output_tag    = "phi_verify"
\end{lstlisting}
\end{tcolorbox}

%------------------------------------------------------------------------------
\subsection*{E27 --- Full Surface-Analysis Showcase}
%------------------------------------------------------------------------------

\textbf{Goal:} capstone --- five probes, all metrics, species filters, and
IVec coupling on a NaCl\,+\,Ar system.\\
\textbf{Key concepts:} all surface metrics simultaneously \textbullet{}
five named probes \textbullet{} integration smoke-test

\begin{tcolorbox}[scriptbox, scriptheader,
  title={scripts/e27\_surface\_full\_analysis\_showcase.vsim (excerpt)}]
\begin{lstlisting}[style=vsimlst]
# surf_A: all metrics, rectangle
[[object.surface]]
name                  = "surf_A_full"
geometry              = "rectangle"
normal_z              = 1.0
width                 = 30.0
height                = 30.0
log_crossings         = true
compute_flux          = true
compute_energy_flux   = true
compute_momentum_flux = true
compute_mass_flux     = true
output_tag            = "surf_A"

# surf_E: IVec flux at crystal top
[[object.surface]]
name              = "surf_E_ivec_top"
geometry          = "disk"
center_z          = 15.0
normal_z          = 1.0
radius            = 8.0
log_crossings     = true
compute_ivec_flux = true
output_tag        = "surf_E_ivec"
\end{lstlisting}
\end{tcolorbox}

\textbf{Sidecar layout} (\texttt{e27.surface.json}):

\begin{tcolorbox}[enhanced, breakable,
  colback=vsimcodebg, colframe=vsimrule, arc=3pt, boxrule=0.6pt,
  left=4pt, right=4pt, top=3pt, bottom=3pt]
\begin{lstlisting}[style=jsonlst]
{
  "run": "e27_surface_full_showcase",
  "surfaces": {
    "surf_A":    { "crossings":[...], "flux":..., "energy_flux":...,
                   "momentum_flux":..., "mass_flux":... },
    "surf_B_Na": { "crossings":[...], "species_flux":{"Na":...} },
    "surf_C_Cl": { "crossings":[...], "species_flux":{"Cl":...} },
    "surf_D_Ar": { "crossings":[...], "mass_flux":... },
    "surf_E_ivec":{ "crossings":[...], "ivec_flux":[...] }
  }
}
\end{lstlisting}
\end{tcolorbox}

%------------------------------------------------------------------------------
\subsection*{E28 --- Full Export Suite (All Formats + Surface)}
%------------------------------------------------------------------------------

\textbf{Goal:} activate every export format simultaneously alongside surface
analysis to verify complete pipeline integration.\\
\textbf{Key concepts:} \texttt{write\_xyz}, \texttt{write\_xyz\_full},
\texttt{write\_analysis\_json}, \texttt{write\_report\_md}, \texttt{write\_rdf}
\textbullet{} surface sidecar coexists with all formats
\textbullet{} integration smoke-test

\begin{tcolorbox}[scriptbox, scriptheader,
  title={scripts/e28\_surface\_export\_suite.vsim (excerpt)}]
\begin{lstlisting}[style=vsimlst]
[analysis.structure]
enabled     = true
compute_rdf = true
rdf_rmax    = 10.0
rdf_nbins   = 200

[analysis.ivec]
enabled         = true
write_json      = true
convergence_eps = 1.0e-4

[export]
write_xyz           = true
write_xyz_full      = true
write_analysis_json = true
write_report_md     = true
write_rdf           = true
output_dir          = "out/e28"

[object.surface]
enabled    = true
write_json = true
output_dir = "out/e28/surfaces"

[[object.surface]]
name                = "export_suite_probe"
geometry            = "disk"
normal_z            = 1.0
radius              = 15.0
log_crossings       = true
compute_flux        = true
compute_energy_flux = true
compute_ivec_flux   = true
output_tag          = "export_suite"
\end{lstlisting}
\end{tcolorbox}

\paragraph{Output files produced by E28}

\begin{center}\footnotesize
\begin{tabular}{@{}ll@{}}
\toprule
File & Source \\
\midrule
\texttt{e28.xyz} & \texttt{write\_xyz = true} \\
\texttt{e28.xyzFull} & \texttt{write\_xyz\_full = true} \\
\texttt{e28.analysis.json} & \texttt{write\_analysis\_json = true} \\
\texttt{e28.report.md} & \texttt{write\_report\_md = true} \\
\texttt{e28.rdf.json} & \texttt{write\_rdf = true} \\
\texttt{e28.identity.json} & \texttt{[analysis.ivec] write\_json = true} \\
\texttt{e28.surface.json} & \texttt{[object.surface] write\_json = true} \\
\bottomrule
\end{tabular}
\end{center}

%------------------------------------------------------------------------------
\subsection*{E29 --- Surface Error Analysis and Diagnostic Patterns}
%------------------------------------------------------------------------------

\textbf{Goal:} intentionally misconfigure surface probes to demonstrate error
handling, validation warnings, and diagnostic output.\\
\textbf{Key concepts:} parser robustness \textbullet{} validation warnings
(non-fatal) \textbullet{} diagnostic patterns \textbullet{} failure-mode documentation

\begin{tcolorbox}[scriptbox, scriptheader,
  title={scripts/e29\_surface\_error\_analysis.vsim (excerpt)}]
\begin{lstlisting}[style=vsimlst]
[object.surface]
enabled    = true
write_json = true
output_dir = "out/e29/surfaces"

# Probe 1: disk with ignored width/height fields
[[object.surface]]
name               = "disk_with_excess_fields"
geometry           = "disk"
radius             = 10.0
width              = 999.0  # Ignored for disk geometry
height             = 999.0  # Ignored for disk geometry
log_crossings      = true
compute_flux       = true
output_tag         = "disk_excess"

# Probe 2: species_filter referencing nonexistent species
[[object.surface]]
name              = "unknown_species_filter"
geometry          = "rectangle"
normal_y          = 1.0
width             = 20.0
height            = 20.0
log_crossings     = true
compute_flux      = true
species_filter    = ["Nonexistent", "Ar"]
output_tag        = "unknown_species"

# Probe 3: compute_ivec_flux without [analysis.ivec]
[[object.surface]]
name              = "orphan_ivec_flux"
geometry          = "rectangle"
normal_z          = 1.0
width             = 15.0
height            = 15.0
log_crossings     = false
compute_ivec_flux = true
output_tag        = ""  # Empty tag -> fallback to "surface_<index>"
\end{lstlisting}
\end{tcolorbox}

\paragraph{Expected warnings (non-fatal)}

\begin{itemize}[leftmargin=1em,itemsep=0pt]
\item \texttt{"species\_filter includes unknown species: Nonexistent"}
\item \texttt{"compute\_ivec\_flux requires [analysis.ivec] enabled; ignored"}
\item \texttt{"output\_tag empty; using fallback `surface\_2'"}
\end{itemize}

%==============================================================================
\newpage
\subsection*{QR --- Updated Section \& Field Reference}
%==============================================================================

\begin{center}\footnotesize
\begin{tabular}{@{}lll@{}}
\toprule
Section & Purpose & WO \\
\midrule
\texttt{[analysis.structure]}        & RDF, coordination number         & --- \\
\texttt{[analysis.sampling]}         & Scalar sampling, MSD             & --- \\
\texttt{[analysis.scale\_sampling]}  & Field projection, RVE            & WO-61D \\
\texttt{[analysis.inference]}        & Property inference               & WO-61B \\
\texttt{[analysis.ikk\_end\_tag]}    & IKK report end-tag badge         & WO-75A \\
\texttt{[analysis.ivec]}             & I-vector $\to$ \texttt{.identity.json} & WO-75B \\
\texttt{[object.<layer>.<basis>]}    & MCF-CAI $3\times3$ state grid    & WO-76 \\
\texttt{[object.surface]}            & Surface analysis master block    & Surface \\
\texttt{[[object.surface]]}          & Per-surface probe definition     & Surface \\
\bottomrule
\end{tabular}
\end{center}

\medskip
\textbf{\texttt{[[object.surface]]} field quick-reference:}

\begin{center}\footnotesize
\begin{tabular}{@{}llll@{}}
\toprule
Field & Type & Default & Notes \\
\midrule
\texttt{name}          & string & \texttt{""} & Probe label \\
\texttt{geometry}      & string & \texttt{"rectangle"} & \texttt{"rectangle"} $|$ \texttt{"disk"} $|$ \texttt{"sphere"} \\
\texttt{center\_x/y/z} & float  & \texttt{0.0} & Centre in \AA \\
\texttt{normal\_x/y/z} & float  & \texttt{0,0,1} & Unit normal \\
\texttt{width}/\texttt{height} & float & \texttt{10.0} & Rectangle half-extents in \AA \\
\texttt{radius}        & float  & \texttt{5.0} & Disk/sphere radius in \AA \\
\texttt{log\_crossings}       & bool & \texttt{false} & Record every event \\
\texttt{compute\_flux}        & bool & \texttt{false} & $\Phi$ metric \\
\texttt{compute\_mass\_flux}  & bool & \texttt{false} & $J_m$ metric \\
\texttt{compute\_energy\_flux} & bool & \texttt{false} & $Q$ metric \\
\texttt{compute\_momentum\_flux} & bool & \texttt{false} & $P$ metric \\
\texttt{compute\_species\_flux} & bool & \texttt{false} & $J_s$ per-species table \\
\texttt{compute\_ivec\_flux}  & bool & \texttt{false} & $J_{iv}$ (needs \texttt{analysis.ivec}) \\
\texttt{species\_filter}      & list & \texttt{[]} & Restrict probe by species \\
\texttt{output\_tag}          & string & \texttt{""} & Key in \texttt{.surface.json} \\
\bottomrule
\end{tabular}
\end{center}

\vfill
\hrule
\smallskip
{\footnotesize\color{vsimgrey}
See also: \texttt{VSIM\_REFERENCE.md} --- complete field reference\quad|\quad
\texttt{docs/VSIM\_LANGUAGE.md} --- language format and registry internals\quad|\quad
\texttt{VSIM\_DEVELOPMENT.md} --- how to add new sections}

\end{document}
